\section{Introduction}

Distributed Asynchronous Object Storage (DAOS) is a high-performance, open-source parallel I/O stack built on Non-Volatile Memory Express (NVMe) and Remote Direct Memory Access (RDMA) technologies \cite{daos-exascale-storage2016sc, daos-intro2020, daos-dram-intro2023, daos-repo}. Its underlying design of shifting from POSIX-based filesystems to an object-based key-value store makes it a promising candidate for next-generation distributed storage systems, especially in high-performance computing (HPC) and large language model (LLM) environments \cite{hdf5-daos2022tpds, daos-adios2024ISC, daos2024weather_flow, mei2024sc-poster, daos-smartnic2025}.

\textcolor{blue}{
In the latest IO500 \cite{io500-ics26} Production list (ICS26, June/24 2026)\dots \cite{io500-ics26}
}


\textcolor{blue}{
In this paper, we present a comprehensive performance benchmark of DAOS, evaluating its I/O throughput, latency, and scalability under various workloads and configurations. We compare DAOS with traditional storage systems to highlight its advantages and identify potential bottlenecks. Our results demonstrate that DAOS can significantly improve I/O performance in HPC and LLM applications, making it a compelling choice for future storage solutions.
}
